\clearpage\fancyhead[L]{\small COMET.}\hypertarget{native-343-lower}{}\label{native-343-lower}
\begin{center}\textbf{COMET.}\par\vspace{4pt}\hrule\vspace{4pt}\fontsize{10}{12}\selectfont Monatliche Correspondenz. Band XXIV. S. 406\,..\,411. 1811 October.\par\vspace{4pt}\hrule\vspace{4pt}\end{center}

\par\indent Die im vorigen Heft geäusserte Hoffnung, dass wir im Stande sein würden, in diesem Stück etwas bestimmtes über die Umlaufszeit des Cometen unsern Lesern mittheilen zu können, ist unerfüllt geblieben. So wahrscheinlich es anfangs war, dass die immer zunehmende starke Abweichung der Beobachtungen von den parabolischen Elementen, in einer reellen Abweichung der Bahn von der Parabel begründet sei, so zeigte es sich doch späterhin, dass eine kleine Correction jener völlig hinreichend war, um eine befriedigende Übereinstimmung beider zu erhalten. Da es nicht ohne Interesse ist zu sehen, wie durch die successiven Änderungen der Elemente, diese zu einer immer bessern Übereinstimmung mit dem Himmel gebracht werden, so fangen wir heute damit an, ein Tableau der Vergleichung sämmtlicher bis zum 11. October gemachten Beobachtungen mit den zum erstenmal verbesserten \textsc{Gauss}ischen Elementen (Mon. Corresp. Sept.-Heft S. 305) [S. 335 d. B.] darzulegen.

\Needspace{41\baselineskip}\SourceLeaf{344}
Vergleichung sämmtlicher Cometen-Beobachtungen mit \textsc{Gauss} parabolischen Elementen.

\begin{center}

\fontsize{10}{12}\selectfont\setlength{\tabcolsep}{7pt}\renewcommand{\arraystretch}{1.05}\begin{tabular}{ll|c|c|l}
\multicolumn{2}{c|}{1811} & \multicolumn{2}{c|}{Abweichung} & Name \\
\multicolumn{2}{c|}{Tag der Beobacht.} & in AR. & in Decl. & des Beobacht. \\\hline
August & 22 & $+61''$ & $-57''$ & Bessel \\
 & 23 & $+70$ & $-79$ & Bessel \\
 & 23 & $+119$ & $-21$ & Olbers \\
 & 25 & $+116$ & $-101$ & Olbers \\
 & 26 & $+98$ & $-84$ & Olbers \\
 & 27 & $+99$ & $-99$ & Olbers \\
 & 27 & $+63$ & $-50$ & Bessel \\
 & 28 & $+34$ & $-68$ & Bessel \\
 & 28 & $+46$ & $-124$ & Olbers \\
 & 29 & $+128$ & $-26$ & Oriani \\
 & 30 & $+9$ & $-41$ & Olbers \\
 & 30 & $+43$ & $-55$ & Bessel \\
 & 31 & $+37$ & $-49$ & Oriani \\
 & 31 & $+47$ & $-120$ & Olbers \\
Septbr. & 1 & $+25$ & $-43$ & Oriani \\
 & 2 & $+19$ & $-74$ & Oriani \\
 & 3 & $+10$ & $-92$ & Oriani \\
 & 3 & $+23$ & $-67$ & Olbers \\
 & 4 & $+25$ & $-73$ & Oriani \\
 & 4 & $+51$ & $-129$ & Olbers \\
 & 4 & $+30$ & $-9$ & Gauss \\
 & 6 & $-43$ & $-89$ & Oriani \\
 & 6 & $-7$ & $-50$ & Lindenau \\
 & 6 & $+9$ & $-11$ & Gauss \\
 & 7 & $-26$ & $-74$ & Lindenau \\
 & 7 & $-17$ & $-80$ & Bessel \\
 & 7 & $-54$ & $+20$ & Gauss \\
 & 8 & $-54$ & $-45$ & Oriani \\
 & 8 & $-41$ & $-44$ & Lindenau \\
 & 9 & $-61$ & $-38$ & Oriani \\
 & 9 & $-62$ & $-7$ & Lindenau \\
 & 9 & $-65$ & $-64$ & Gauss \\
\end{tabular}
\end{center}
\Needspace{38\baselineskip}
\begin{center}

\fontsize{10}{12}\selectfont\setlength{\tabcolsep}{7pt}\renewcommand{\arraystretch}{1.05}\begin{tabular}{ll|c|c|l}
\multicolumn{2}{c|}{1811} & \multicolumn{2}{c|}{Abweichung} & Name \\
\multicolumn{2}{c|}{Tag der Beobacht.} & in AR. & in Decl. & des Beobacht. \\\hline
Septbr. & 10 & $-116''$ & $-27$ & Lindenau \\
 & 10 & $-54$ & $-71$ & Schubert \\
 & 11 & $-70$ & $-73$ & Oriani \\
 & 11 & $-135$ & $-28$ & Lindenau \\
 & 11 & $-82$ & $-106$ & Bessel. \\
 & 12 & $-120$ & $-43$ & Oriani \\
 & 12 & $-37$ & $-78$ & Bessel \\
 & 12 & $-86$ & $-53$ & Schubert \\
 & 13 & $-111$ & $-118$ & Oriani \\
 & 13 & $-130$ & $-34$ & Lindenau \\
 & 13 & $-119$ & $-74$ & Schubert \\
 & 14 & $-135$ & $-94$ & Oriani \\
 & 14 & $-146$ & $-34$ & Lindenau \\
 & 14 & $-140$ & .\;.\;. & Gauss \\
 & 14 & $-139$ & $-30$ & Schubert \\
 & 15 & $-154$ & $-51$ & Lindenau \\
 & 15 & $-195$ & $-76$ & Gauss \\
 & 15 & $-158$ & $-65$ & Schubert \\
 & 16 & $-225$ & $-54$ & Lindenau \\
 & 16 & $-219$ & $-30$ & Gauss \\
 & 16 & $-184$ & $-52$ & Schubert \\
 & 17 & $-172$ & $-43$ & Oriani \\
 & 17 & $-261$ & $-29$ & Lindenau \\
 & 18 & $-247$ & $-34$ & Lindenau \\
 & 19 & $-287$ & $-66$ & Lindenau \\
 & 22 & $-443$ & .\;.\;. & Bessel \\
 & 23 & $-420$ & .\;.\;. & Bessel \\
 & 26 & $-626$ & $+53$ & Lindenau \\
 & 30 & $-839$ & $+130$ & Lindenau \\
Octob. & 2 & $-879$ & $+118$ & Gauss \\
 & 5 & $-1123$ & $+336$ & Lindenau \\
 & 11 & $-1271$ & $+626$ & Lindenau \\
\end{tabular}

\end{center}

\par\indent Herr Professor \textsc{Gauss}, der uns in den letzten Tagen des Septembers und Anfang October das grosse Vergnügen gewährte, einen 14tägigen Aufenthalt auf der hiesigen Sternwarte zu machen, war sehr geneigt, in diesen Abweichungen die Spur einer elliptischen oder hyperbolischen Laufbahn zu sehen, und fing schon hier einige auf deren Bestimmung Bezug habende Rechnungen an, ohne jedoch zu einem bestimmten Resultat darüber zu gelangen. Allein schon am zweiten Tage nach seiner Rückkunft in Göttingen, theilte er uns seine erhaltenen interessanten Resultate mit, die wir hier mit dessen eigenen Worten folgen lassen: „d. d. Göttingen, am 14. October. Meine letzten Rechnungen über den Cometen habe ich sofort nachgesehen, und wie ich voraus vermuthete, bald einen Schreibfehler von 10 Minuten darinnen entdeckt, nach dessen Verbesserung das Resultat ganz anders und zwar dahin ausgefallen ist,

\begin{center}
\emph{dass von einer Ellipticität der Bahn noch gar keine sichere Spur zu bemerken ist.}
\end{center}

\noindent Meine verbesserten parabolischen Elemente sind folgende:

\Needspace{12\baselineskip}
\begin{center}\fontsize{10}{12}\selectfont\setlength{\tabcolsep}{4pt}\renewcommand{\arraystretch}{1.15}\begin{tabular}{l}
Durchgang durch das Perihel 1811 Sept. 12\\
\quad $5^{\mathrm u}\,32^{\mathrm m}\,46^{\mathrm s}$ M. Z. in Göttingen\\
\end{tabular}\end{center}
\begin{center}\fontsize{10}{12}\selectfont\setlength{\tabcolsep}{3pt}\renewcommand{\arraystretch}{1.15}\begin{tabular}{llr}
Länge des Perihels & .\;.\;.\;.\;.\;.\;. & $75^\circ\,4'\,43''$\\
Log. des Abstandes im Perihel & .\;.\;. & $0.015530$\\
Länge des aufsteigenden Knoten & . & $140^\circ\,21'\,40''$\\
Neigung der Bahn & .\;.\;.\;.\;.\;. & $73\;4\;18$\\
\end{tabular}\end{center}
\noindent Bewegung rückläufig

\Needspace{17\baselineskip}\SourceLeaf{345}
\par\indent Die Constanten für die Coordinaten in Beziehung auf den Aequator finde ich so

\[
\begin{array}{ll}
x = \dfrac{\alpha.\sin(v+348^\circ\,50'\,38'')}{\cos\frac{1}{2}v} & \log \alpha = 9.91435 \\[2ex]
y = \dfrac{\ObservedBeta.\sin(v+171^\circ\,56'\,51'')}{\cos\frac{1}{2}v} & \log \ObservedBeta = 9.80148 \\[2ex]
z = \dfrac{\gamma.\sin(v+80^\circ\;0'\;3'')}{\cos\frac{1}{2}v} & \log \gamma = 0.01538 \\
\end{array}
\]

\par\indent Diese Elemente sind übrigens nur flüchtig berechnet und würden sich leicht den Beobachtungen noch besser anpassen lassen; ich hielt es aber nicht der Mühe werth, dies schon jetzt zu thun, da sie hinreichen, die Beobachtungen noch einige Zeit damit zu vergleichen, und die bis jetzt Statt findende Unmöglichkeit, etwas sicheres über die Ellipticität der Bahn zu sagen, zeigen. Ich glaube nicht, dass die Umlaufszeit, wenn die Bahn elliptisch ist, unter 1000 Jahr sein kann. Ich habe so eben diese Elemente mit Ihren und meinen Beobachtungen verglichen, wodurch das, was ich vorhin sagte bestätigt wird.

\Needspace{25\baselineskip}
\begin{center}\fontsize{10}{12}\selectfont\setlength{\tabcolsep}{6pt}\renewcommand{\arraystretch}{1.15}\begin{tabular}{l|c|c|c|c}
 & \multicolumn{4}{c}{Unterschied}\\
 & \multicolumn{2}{c|}{in gerader Aufst.} & \multicolumn{2}{c}{in Abweichung} \\
1811 & v. L. & G. & v. L. & G. \\\hline
Septemb. 4 & --- & $+5''$ & --- & $+23''$ \\
6 & --- & $-7$ & --- & $+28$ \\
7 & --- & $-58$ & --- & $+61$ \\
9 & --- & $-38$ & --- & $-18$ \\
10 & $-69''$ & & $+26''$ & \\
11 & $-70$ & & $+39$ & \\
14 & $-12$ & $-9$ & $+26$ & \\
15 & $+5$ & $-17$ & $+7$ & $-18$ \\
16 & $-37$ & $-36$ & $+10$ & $+35$ \\
17 & $-43$ & & $+33$ & \\
18 & $+4$ & & $+16$ & \\
19 & $-27$ & & $+1$ & \\
26 & $-47$ & & $+63$ & \\
30 & --- & $-50$ & $+76$ & \\
October 2 & --- & $-29$ & & $+19$ \\
\end{tabular}
\end{center}

\par\indent Die Elemente geben also --- wie es scheint während der ganzen Zeit dieser Beobachtungen --- die Rectascensionen etwa eine halbe Minute zu klein. Die Declinationen ungefähr eben so viel zu gross, welcher Unterschied sich leicht wegschaffen liesse, ohne die Parabel zu verlassen und ohne die Übereinstimmung mit den von \textsc{Zach}schen Beobachtungen zu verschlechtern. \textsc{Olbers} hat mir in einem hier vorgefundenen Briefe von seiner sinnreichen Hypothese über den Cometen-Schweif eine kurze Nachricht gegeben. Er nimmt einen Stoff an, welcher vom Cometen erzeugt, von diesem und von der Sonne abgestossen, sich da anhäuft, wo beide Repulsionskräfte eine Art von Gleichgewicht halten, und sich in eine Art von hohler parabolischer Conoide formirt. \textsc{Olbers} wird Ihnen darüber ausführlicher schreiben.
